\documentclass[11pt, a4paper]{article}

\usepackage[utf8]{inputenc}
\usepackage[T1]{fontenc}
\usepackage{amsmath, amssymb, amsthm, mathtools}
\usepackage{bm}
\usepackage{booktabs}
\usepackage{graphicx}
\usepackage{hyperref}
\usepackage[margin=2.5cm]{geometry}
\usepackage{enumitem}
\usepackage{xcolor}
\usepackage{tcolorbox}
\usepackage{float}

% Algorithm box (no external package needed)
\newtcolorbox{algorithmbox}[1][]{
    colback=white,
    colframe=black!70,
    fonttitle=\bfseries,
    #1
}

% Theorem environments
\theoremstyle{definition}
\newtheorem{definition}{Definition}[section]
\newtheorem{example}{Example}[section]
\theoremstyle{plain}
\newtheorem{theorem}{Theorem}[section]
\newtheorem{proposition}{Proposition}[section]
\newtheorem{corollary}{Corollary}[section]
\newtheorem{lemma}{Lemma}[section]
\theoremstyle{remark}
\newtheorem{remark}{Remark}[section]

% Colored boxes
\newtcolorbox{keyresult}[1][]{
    colback=blue!5!white,
    colframe=blue!60!black,
    fonttitle=\bfseries,
    title={Key Insight},
    #1
}
\newtcolorbox{warningbox}[1][]{
    colback=red!5!white,
    colframe=red!60!black,
    fonttitle=\bfseries,
    title={Warning},
    #1
}
\newtcolorbox{refbox}[1][]{
    colback=gray!5!white,
    colframe=gray!50!black,
    fonttitle=\bfseries,
    title={Go Deeper},
    #1
}
\newtcolorbox{keyidea}[1][]{
    colback=green!5!white,
    colframe=green!50!black,
    fonttitle=\bfseries,
    title={Core Idea},
    #1
}
\newtcolorbox{prereq}[1][]{
    colback=orange!5!white,
    colframe=orange!60!black,
    fonttitle=\bfseries,
    title={Read Before Coding},
    #1
}
\newtcolorbox{objective}[1][]{
    colback=blue!5!white,
    colframe=blue!60!black,
    fonttitle=\bfseries,
    title={What You Will Build},
    #1
}

% Custom commands
\newcommand{\VaR}{\operatorname{VaR}}
\newcommand{\CVaR}{\operatorname{CVaR}}
\newcommand{\ES}{\operatorname{ES}}
\newcommand{\EV}{\mathbb{E}}
\newcommand{\PV}{\mathbb{P}}
\newcommand{\RV}{\mathbb{R}}
\newcommand{\ind}{\boldsymbol{1}}

\title{%
    \textbf{Reinforcement Learning for Portfolio Risk Management} \\[0.5em]
    \large CVaR-Constrained Dynamic Allocation
}
\author{Andr\'es Gonz\'alez Ortega}
\date{March 2026}

\begin{document}
\maketitle

\begin{abstract}
An applied treatment of reinforcement learning for dynamic portfolio allocation
under tail-risk constraints. We formulate the allocation problem as a Constrained
Markov Decision Process (CMDP) where the agent maximises risk-adjusted returns
subject to a Conditional Value-at-Risk (CVaR) budget. Two optimisation paradigms
are developed: policy-gradient (REINFORCE) and derivative-free Natural Evolution
Strategies (NES). Reward shaping incorporates rolling CVaR penalties and
transaction costs to produce a realistic trading signal. The RL agent is
benchmarked against equal-weight, mean-variance, and risk-parity strategies,
with the goal of achieving competitive returns while maintaining tighter control
over drawdowns and tail risk. The document provides the mathematical foundations
required to implement and interpret the code in Project~13.
\end{abstract}

\tableofcontents
\newpage

% ============================================================
\section{Portfolio Allocation as a Sequential Decision}
\label{sec:sequential}
% ============================================================

\begin{prereq}
This section assumes familiarity with basic portfolio theory (mean-variance
optimisation) and Markov Decision Processes. We use CVaR as the central
risk measure; see Project~05 (EVT \& Tail Risk) for a self-contained treatment.
\end{prereq}

\subsection{The Static Paradigm and Its Limits}

The classical framework of Markowitz~(1952) solves a \emph{single-period}
problem:
\begin{equation}
    \max_{\bm{w}} \; \bm{w}^\top \bm{\mu}
    - \frac{\gamma}{2}\,\bm{w}^\top \bm{\Sigma}\,\bm{w},
    \qquad
    \text{s.t.}\quad \bm{1}^\top \bm{w} = 1, \quad w_i \geq 0,
    \label{eq:markowitz}
\end{equation}
where $\bm{w} \in \RV^N$ are portfolio weights, $\bm{\mu}$ is the vector of
expected returns, $\bm{\Sigma}$ is the covariance matrix, and $\gamma > 0$ is a
risk-aversion parameter. The solution is computed once and held until the next
rebalancing date.

This approach has two fundamental problems:
\begin{enumerate}[label=(\roman*)]
    \item \textbf{Stationarity assumption.} The inputs $(\bm{\mu}, \bm{\Sigma})$
          are treated as fixed, yet financial markets exhibit regime changes,
          time-varying volatility, and momentum effects.
    \item \textbf{No adaptation.} A portfolio optimised for calm markets becomes
          dangerously concentrated when correlations spike during crises.
\end{enumerate}

\begin{keyidea}
The core insight motivating this project: portfolio allocation is not a
single optimisation problem but a \emph{sequential decision process} where
the agent must adapt to changing market conditions at every rebalancing step.
\end{keyidea}

\subsection{MDP Formulation}

We model portfolio rebalancing as a discrete-time Markov Decision Process
$(\mathcal{S}, \mathcal{A}, P, R, \gamma_d)$:

\begin{definition}[Portfolio MDP]
\label{def:portfolio_mdp}
The components are:
\begin{itemize}[leftmargin=2em]
    \item \textbf{State} $s_t \in \mathcal{S}$: a feature vector encoding the
          agent's information at time $t$,
          \[
              s_t = \bigl(\bm{r}_{t-L:t},\; \bm{w}_{t}^{-},\;
              \sigma_{t}^{\text{rolling}},\; \text{DD}_t\bigr),
          \]
          where $\bm{r}_{t-L:t}$ are the most recent $L$ asset returns,
          $\bm{w}_t^{-}$ are the current (pre-rebalance) portfolio weights,
          $\sigma_t^{\text{rolling}}$ is a rolling volatility estimate, and
          $\text{DD}_t$ is the current drawdown from peak.

    \item \textbf{Action} $a_t \in \mathcal{A}$: the new portfolio weight
          vector $\bm{w}_t^{+}$ satisfying $\bm{1}^\top \bm{w}_t^{+} = 1$
          and $w_{t,i}^{+} \geq 0$ (long-only).

    \item \textbf{Transition} $P(s_{t+1} \mid s_t, a_t)$: determined by
          market dynamics (asset returns), which the agent does not control.

    \item \textbf{Reward} $r_t = R(s_t, a_t, s_{t+1})$: a risk-adjusted
          portfolio return (detailed in Section~\ref{sec:reward}).

    \item \textbf{Discount factor} $\gamma_d \in (0, 1]$: controls the
          agent's planning horizon.
\end{itemize}
\end{definition}

The agent learns a \emph{policy} $\pi_\theta(a \mid s)$ parameterised by
$\theta$ that maps observed market conditions to portfolio allocations. The
objective is:
\begin{equation}
    \max_\theta \; J(\theta) = \EV_{\pi_\theta}
    \biggl[\sum_{t=0}^{T-1} \gamma_d^t\, r_t \biggr].
    \label{eq:rl_objective}
\end{equation}

\begin{remark}
Unlike supervised learning, there is no ``correct'' label for the optimal
allocation. The agent discovers good strategies through \emph{interaction}
with the environment (historical or simulated market data), guided only by
the reward signal.
\end{remark}


% ============================================================
\section{CVaR-Constrained MDP}
\label{sec:cvar_cmdp}
% ============================================================

\subsection{Why Constrain Tail Risk?}

Maximising expected return (even risk-adjusted return) can still produce
unacceptable tail losses. A portfolio with a high Sharpe ratio may
simultaneously have catastrophic drawdowns during rare events. We want a
\emph{hard constraint} on how bad the worst outcomes can be.

\begin{definition}[CVaR at level $\alpha$]
\label{def:cvar}
For a loss random variable $L$ and confidence level $\alpha \in (0,1)$,
\[
    \VaR_\alpha(L) = \inf\{l : \PV(L \leq l) \geq \alpha\},
\]
\[
    \CVaR_\alpha(L)
    = \frac{1}{1-\alpha}\int_\alpha^1 \VaR_u(L)\,du
    = \EV\bigl[L \mid L \geq \VaR_\alpha(L)\bigr].
\]
CVaR is the expected loss in the worst $(1-\alpha)$ fraction of scenarios.
It is a \emph{coherent} risk measure: subadditive, positive homogeneous,
translation invariant, and monotone.
\end{definition}

\subsection{The Constrained Optimisation Problem}

We augment the RL objective~\eqref{eq:rl_objective} with a CVaR constraint:
\begin{equation}
    \max_\theta \; J(\theta)
    \qquad \text{subject to} \qquad
    \CVaR_\alpha\!\bigl(L^{\pi_\theta}\bigr) \leq \bar{c},
    \label{eq:cmdp}
\end{equation}
where $L^{\pi_\theta}$ is the portfolio loss distribution induced by policy
$\pi_\theta$ and $\bar{c}$ is the CVaR budget (threshold).

\subsection{Lagrangian Relaxation}

Direct enforcement of~\eqref{eq:cmdp} is intractable in the RL loop. We
relax it via a Lagrangian:
\begin{equation}
    \mathcal{L}(\theta, \lambda)
    = J(\theta) - \lambda\,
      \max\!\bigl(\CVaR_\alpha(L^{\pi_\theta}) - \bar{c},\; 0\bigr),
    \label{eq:lagrangian}
\end{equation}
where $\lambda \geq 0$ is the Lagrange multiplier.

\begin{keyresult}
The multiplier $\lambda$ is the \textbf{shadow price of risk}. It encodes
how aggressively the agent trades return for tail-risk reduction:
\begin{itemize}[leftmargin=2em]
    \item When $\CVaR_\alpha < \bar{c}$: the constraint is \emph{inactive},
          and the agent maximises return freely ($\lambda$ contribution is zero).
    \item When $\CVaR_\alpha > \bar{c}$: the penalty activates, pushing the
          agent toward more conservative allocations.
\end{itemize}
A higher $\lambda$ produces a more risk-averse agent; a lower $\lambda$
tolerates larger tail losses in exchange for higher expected returns.
\end{keyresult}

\subsection{Dual Ascent on $\lambda$}

In practice $\lambda$ can be treated as a fixed hyperparameter or updated
via dual gradient ascent:
\begin{equation}
    \lambda_{k+1}
    = \max\!\bigl(0,\;
      \lambda_k + \eta_\lambda\,
      (\widehat{\CVaR}_\alpha - \bar{c})\bigr),
    \label{eq:dual_ascent}
\end{equation}
where $\widehat{\CVaR}_\alpha$ is the empirical CVaR estimated from
recent rollouts and $\eta_\lambda$ is the dual learning rate. This drives
$\lambda$ upward when the constraint is violated and downward when there
is slack.

\begin{warningbox}
Dual ascent introduces an \emph{inner--outer} loop: the inner loop updates
$\theta$ (policy parameters) for a fixed $\lambda$; the outer loop updates
$\lambda$ based on whether the CVaR constraint is satisfied. Convergence
requires the inner loop to approximately solve its sub-problem before each
$\lambda$ update. In practice, updating $\lambda$ every $K$ policy steps
(with $K$ between 5 and 20) works well.
\end{warningbox}


% ============================================================
\section{Policy Gradient and Evolutionary Strategies}
\label{sec:optimisation}
% ============================================================

We consider two families of algorithms to optimise $\theta$.

\subsection{REINFORCE (Policy Gradient)}

Given a parameterised stochastic policy $\pi_\theta(a \mid s)$, the
policy gradient theorem (Williams, 1992) gives:
\begin{equation}
    \nabla_\theta J(\theta)
    = \EV_{\pi_\theta}\!\biggl[
        \sum_{t=0}^{T-1} \nabla_\theta \log \pi_\theta(a_t \mid s_t)\,
        \hat{A}_t
      \biggr],
    \label{eq:reinforce}
\end{equation}
where $\hat{A}_t$ is an advantage estimate (e.g.\ return-to-go minus a
baseline). The gradient is estimated by sampling trajectories and computing
the log-probability-weighted returns.

\begin{algorithmbox}[title={REINFORCE for Portfolio Allocation}]
\begin{enumerate}[leftmargin=1.5em]
    \item Sample $M$ episodes $\{\tau^{(m)}\}_{m=1}^M$ under $\pi_\theta$.
    \item For each episode, compute returns
          $G_t^{(m)} = \sum_{k=t}^{T-1}\gamma_d^{k-t} r_k^{(m)}$.
    \item Estimate gradient:
          $\hat{g} = \frac{1}{M}\sum_{m=1}^M \sum_t
          \nabla_\theta \log\pi_\theta(a_t^{(m)} \mid s_t^{(m)})\,
          (G_t^{(m)} - b)$.
    \item Update: $\theta \leftarrow \theta + \eta\,\hat{g}$.
\end{enumerate}
\end{algorithmbox}

\textbf{Pros:} Principled gradient estimate; works with any differentiable
policy (neural network with softmax output for weights).

\textbf{Cons:} High variance; requires careful baseline tuning; the CVaR
penalty involves sorting (non-differentiable), which complicates gradient
computation.

\subsection{Natural Evolution Strategies (NES)}

NES (Salimans et al., 2017) is a \emph{derivative-free} black-box
optimiser. Instead of computing gradients through the policy, it
maintains a search distribution over parameters $\theta \sim
\mathcal{N}(\bm{\mu}, \sigma^2 I)$ and uses fitness evaluations to
estimate the natural gradient.

\begin{definition}[NES Update]
\label{def:nes}
Given a population of $N_p$ perturbations
$\{\bm{\epsilon}_i\}_{i=1}^{N_p}$,
$\bm{\epsilon}_i \sim \mathcal{N}(\bm{0}, I)$, evaluate the fitness
$F(\bm{\mu} + \sigma\,\bm{\epsilon}_i)$ for each perturbation. The
parameter update is:
\begin{equation}
    \bm{\mu} \leftarrow \bm{\mu}
    + \frac{\eta}{N_p\,\sigma}
      \sum_{i=1}^{N_p} F_i\,\bm{\epsilon}_i,
    \label{eq:nes_update}
\end{equation}
where $F_i = F(\bm{\mu} + \sigma\,\bm{\epsilon}_i)$ is the fitness of
the $i$-th perturbation.
\end{definition}

\textbf{Mirror sampling.} For each $\bm{\epsilon}_i$, also evaluate
$-\bm{\epsilon}_i$, producing an antithetic pair. This halves the
variance of the gradient estimate at the cost of doubling evaluations
(a worthwhile trade in practice).

\begin{algorithmbox}[title={NES with Mirror Sampling for CVaR-Constrained Allocation}]
\begin{enumerate}[leftmargin=1.5em]
    \item Initialise $\bm{\mu}$, perturbation scale $\sigma$,
          learning rate $\eta$, CVaR multiplier $\lambda$.
    \item For each generation $k = 1, 2, \ldots$:
    \begin{enumerate}[label=(\alph*)]
        \item Sample $N_p/2$ perturbations $\bm{\epsilon}_i \sim
              \mathcal{N}(\bm{0}, I)$.
        \item Evaluate fitness for $\bm{\mu} + \sigma\bm{\epsilon}_i$
              and $\bm{\mu} - \sigma\bm{\epsilon}_i$
              (mirror sampling).
        \item Fitness = Lagrangian~\eqref{eq:lagrangian}: episode return
              minus $\lambda \cdot \max(\widehat{\CVaR} - \bar{c}, 0)$.
        \item Update $\bm{\mu}$ via~\eqref{eq:nes_update}.
        \item Optionally update $\lambda$ via~\eqref{eq:dual_ascent}.
    \end{enumerate}
\end{enumerate}
\end{algorithmbox}

\begin{keyidea}
\textbf{Why NES for risk-constrained allocation?}
\begin{enumerate}[leftmargin=2em]
    \item \textbf{Non-differentiable objectives.} CVaR computation
          involves sorting returns and taking the tail mean --- operations
          that break standard backpropagation. NES treats the objective as
          a black box.
    \item \textbf{Robustness to noisy rewards.} Financial returns are
          inherently noisy. NES aggregates over a population of
          perturbations, smoothing out noise.
    \item \textbf{Parallelisable.} Each perturbation can be evaluated
          independently, making NES well-suited for GPU or multi-core
          execution.
\end{enumerate}
This is the same optimisation strategy used in Project~08 (Deep Hedging)
for derivative-free policy search.
\end{keyidea}


% ============================================================
\section{Reward Shaping for Risk}
\label{sec:reward}
% ============================================================

The reward signal is the mechanism through which we communicate our
risk--return preferences to the agent. A poorly designed reward yields
an agent that maximises the wrong quantity.

\subsection{Composite Reward}

At each rebalancing step $t$, the reward is:
\begin{equation}
    r_t = \underbrace{R_t^{\text{port}}}_{\text{portfolio return}}
    \;-\; \underbrace{\lambda\,
    \max\!\bigl(\widehat{\CVaR}_{\alpha,t} - \bar{c},\; 0\bigr)
    }_{\text{CVaR penalty}}
    \;-\; \underbrace{\kappa\,
    \sum_{i=1}^{N}\bigl|w_{t,i}^{+} - w_{t,i}^{-}\bigr|
    \cdot c_{\text{rate}}
    }_{\text{transaction cost}},
    \label{eq:reward}
\end{equation}
where:
\begin{itemize}[leftmargin=2em]
    \item $R_t^{\text{port}} = \bm{w}_t^{+\top}\bm{r}_t$ is the
          realised one-step portfolio return.
    \item $\widehat{\CVaR}_{\alpha,t}$ is the empirical CVaR computed
          from a rolling window of recent portfolio returns (e.g.\ the
          past 60 trading days).
    \item $\lambda \geq 0$ is the risk penalty weight (Lagrange multiplier).
    \item $\kappa \geq 0$ scales the importance of trading costs.
    \item $c_{\text{rate}}$ is the per-unit proportional transaction cost
          (e.g.\ 10 basis points).
\end{itemize}

\subsection{Effect of Hyperparameters}

\begin{center}
\begin{tabular}{lll}
\toprule
\textbf{Parameter} & \textbf{Low value} & \textbf{High value} \\
\midrule
$\lambda$ (risk penalty)       & Aggressive, higher returns,
                                 larger drawdowns
                               & Conservative, lower returns,
                                 tighter tail risk \\
$\kappa$ (cost penalty)        & Frequent rebalancing,
                                 responsive
                               & Infrequent rebalancing,
                                 stable \\
$\bar{c}$ (CVaR threshold)    & Tight constraint, limits
                                 tail risk early
                               & Loose constraint, allows
                                 more risk \\
$\alpha$ (CVaR confidence)     & Shallow tail ($\alpha=0.90$),
                                 moderate risk
                               & Deep tail ($\alpha=0.99$),
                                 extreme risk \\
\bottomrule
\end{tabular}
\end{center}

\begin{warningbox}
The CVaR penalty and the portfolio return operate on very different scales.
If $\lambda$ is too large relative to typical returns, the agent learns to
hold all-cash or all-bonds. If $\lambda$ is too small, the CVaR constraint
has no effect. In practice, normalise returns and CVaR to similar
magnitudes before tuning $\lambda$.
\end{warningbox}

\subsection{Rolling CVaR Estimation}

Given a window of $W$ recent portfolio returns
$\{R^{\text{port}}_{t-W+1}, \ldots, R^{\text{port}}_t\}$:
\begin{enumerate}[leftmargin=2em]
    \item Convert returns to losses: $l_k = -R^{\text{port}}_k$.
    \item Sort losses in descending order:
          $l_{(1)} \geq l_{(2)} \geq \cdots \geq l_{(W)}$.
    \item The number of tail observations:
          $n_\alpha = \lceil (1-\alpha)\,W \rceil$.
    \item Empirical CVaR:
          $\widehat{\CVaR}_\alpha = \frac{1}{n_\alpha}
          \sum_{k=1}^{n_\alpha} l_{(k)}$.
\end{enumerate}

\begin{remark}
This is a non-parametric estimate. For small windows ($W < 100$) and
high $\alpha$ ($\geq 0.95$), the estimate is noisy because only a few
observations fall in the tail. Using $\alpha = 0.95$ with $W = 60$ yields
approximately $n_\alpha = 3$ tail observations. This noise is part of
why derivative-free optimisation (NES) is preferred over gradient methods.
\end{remark}


% ============================================================
\section{Benchmark Strategies}
\label{sec:benchmarks}
% ============================================================

To evaluate the RL agent, we compare it against three classical strategies
that span the spectrum from estimation-free to model-dependent.

\subsection{Equal Weight (1/N)}

The simplest possible allocation:
\begin{equation}
    w_i = \frac{1}{N}, \qquad i = 1, \ldots, N.
    \label{eq:equal_weight}
\end{equation}

DeMiguel, Garlappi \& Uppal~(2009) showed that $1/N$ outperforms many
``optimal'' strategies in out-of-sample tests. The reason is instructive:
optimal strategies introduce \emph{estimation error} through
$(\hat{\bm{\mu}}, \hat{\bm{\Sigma}})$, and this error often exceeds the
benefit of optimisation.

\begin{keyresult}
$1/N$ is the natural baseline: any strategy that cannot beat equal weighting
is not worth its complexity. Maximum diversification with zero estimation
error.
\end{keyresult}

\subsection{Mean-Variance (Markowitz)}

The tangency portfolio weights (assuming a risk-free rate of zero):
\begin{equation}
    \bm{w}^{\text{MV}}
    = \frac{\bm{\Sigma}^{-1}\bm{\mu}}
           {\bm{1}^\top\bm{\Sigma}^{-1}\bm{\mu}}.
    \label{eq:mv}
\end{equation}

This is optimal under the assumption that returns are i.i.d.\ multivariate
normal and that $(\bm{\mu}, \bm{\Sigma})$ are known exactly. Neither
assumption holds in practice:
\begin{itemize}[leftmargin=2em]
    \item Returns are fat-tailed and exhibit volatility clustering.
    \item Estimated $\hat{\bm{\mu}}$ is notoriously unreliable;
          $\hat{\bm{\Sigma}}$ is more stable but still subject to
          regime changes.
\end{itemize}

\begin{warningbox}
Mean-variance portfolios are highly sensitive to the expected return
vector $\bm{\mu}$. Small perturbations in $\hat{\bm{\mu}}$ can flip
the optimal weights entirely. This is the \emph{Markowitz curse}: the
assets with the highest estimated expected returns also tend to have the
highest estimation error, leading to extreme and unstable allocations.
\end{warningbox}

\subsection{Risk Parity}

Risk parity sets weights so that each asset contributes equally to
portfolio \emph{volatility}:
\begin{equation}
    w_i \propto \frac{1}{\sigma_i},
    \qquad \text{normalised so that } \sum_i w_i = 1,
    \label{eq:risk_parity}
\end{equation}
where $\sigma_i = \sqrt{\Sigma_{ii}}$ is the marginal volatility of
asset $i$. The intuition: low-volatility assets receive higher weight,
preventing the portfolio from being dominated by a single volatile asset.

\begin{remark}
Risk parity equalises \emph{volatility} contributions, not
\emph{tail risk} contributions. During crises, correlations spike and
even a volatility-balanced portfolio can suffer concentrated tail losses.
This motivates the tail risk parity approach developed in
Section~\ref{sec:dynamic_risk}.
\end{remark}

\subsection{Summary of Benchmark Properties}

\begin{center}
\begin{tabular}{lccc}
\toprule
\textbf{Strategy} & \textbf{Estimation needed}
                   & \textbf{Adapts to regimes}
                   & \textbf{Controls tail risk} \\
\midrule
Equal weight ($1/N$)  & None     & No  & No  \\
Mean-variance         & $\bm{\mu}, \bm{\Sigma}$ & No & No \\
Risk parity           & $\bm{\Sigma}$ only & Partially & No \\
\textbf{RL agent}     & Learned  & \textbf{Yes} & \textbf{Yes} \\
\bottomrule
\end{tabular}
\end{center}


% ============================================================
\section{Dynamic Risk Budgeting}
\label{sec:dynamic_risk}
% ============================================================

\subsection{Marginal CVaR}

The \emph{marginal CVaR} of asset $i$ measures how sensitive the
portfolio CVaR is to a small change in weight $w_i$:
\begin{equation}
    \text{MCVaR}_i
    = \frac{\partial\,\CVaR_\alpha(\bm{w})}{\partial w_i}
    = \EV\!\bigl[r_i \mid R^{\text{port}} \leq -\VaR_\alpha\bigr],
    \label{eq:mcvar}
\end{equation}
where $R^{\text{port}} = \bm{w}^\top \bm{r}$ and the expectation is
taken over the joint return distribution, conditioned on portfolio
returns being in the worst $(1-\alpha)$ tail.

\begin{definition}[CVaR Contribution]
\label{def:cvar_contribution}
The CVaR contribution of asset $i$ is the product of its weight and
its marginal CVaR:
\[
    \text{CCVaR}_i = w_i \cdot \text{MCVaR}_i.
\]
By Euler's theorem for homogeneous functions, the total portfolio CVaR
decomposes as:
\begin{equation}
    \CVaR_\alpha(\bm{w})
    = \sum_{i=1}^{N} \text{CCVaR}_i
    = \sum_{i=1}^{N} w_i\,\text{MCVaR}_i.
    \label{eq:euler_decomp}
\end{equation}
\end{definition}

\subsection{Tail Risk Parity}

\begin{definition}[Tail Risk Parity]
\label{def:tail_risk_parity}
A portfolio satisfies tail risk parity if every asset contributes
equally to portfolio CVaR:
\begin{equation}
    \text{CCVaR}_1 = \text{CCVaR}_2 = \cdots = \text{CCVaR}_N
    = \frac{\CVaR_\alpha(\bm{w})}{N}.
    \label{eq:tail_parity}
\end{equation}
\end{definition}

While standard risk parity equalises volatility contributions, tail
risk parity equalises \emph{tail} contributions. This is the risk-aware
analogue: it prevents any single asset from dominating the portfolio's
worst-case losses.

\begin{keyidea}
A well-trained RL agent should \textbf{approximately} converge to a tail
risk parity allocation --- not because it is told to, but because the CVaR
penalty in the reward function incentivises equalising tail risk
contributions. If an asset dominates the tail, the agent learns to
underweight it. We validate this by comparing the agent's learned
allocations with the analytical tail risk parity solution.
\end{keyidea}

\subsection{Empirical Marginal CVaR Estimation}

Given $M$ historical scenarios (or simulated paths) of asset returns
$\{\bm{r}^{(m)}\}_{m=1}^M$:
\begin{enumerate}[leftmargin=2em]
    \item Compute portfolio returns:
          $R^{(m)} = \bm{w}^\top \bm{r}^{(m)}$ for each scenario.
    \item Identify the tail set:
          $\mathcal{T} = \{m : R^{(m)} \leq -\VaR_\alpha\}$,
          with $|\mathcal{T}| = \lceil(1-\alpha)M\rceil$.
    \item Estimate marginal CVaR for each asset:
          \[
              \widehat{\text{MCVaR}}_i
              = -\frac{1}{|\mathcal{T}|}
              \sum_{m \in \mathcal{T}} r_i^{(m)}.
          \]
    \item CVaR contribution:
          $\widehat{\text{CCVaR}}_i = w_i \cdot \widehat{\text{MCVaR}}_i$.
\end{enumerate}

\begin{remark}
The RL agent does not compute marginal CVaR explicitly. Instead, it
learns to avoid allocations that concentrate tail risk through the reward
signal. The marginal CVaR analysis is a \emph{diagnostic}: it lets us
check whether the agent's behaviour is consistent with good risk
budgeting.
\end{remark}


% ============================================================
\section{Results Interpretation}
\label{sec:results}
% ============================================================

Four diagnostic plots form the core of the evaluation framework for
the RL agent.

\subsection{Cumulative Return Curves}

\textbf{What it shows.} The growth of \$1 invested under each strategy
over the evaluation period.

\textbf{What to look for.}
\begin{itemize}[leftmargin=2em]
    \item Does the RL agent \emph{outperform} or \emph{match} the
          benchmarks in total return?
    \item Is the growth path \emph{smoother}? A strategy that earns the
          same total return with less volatility is superior.
    \item During \emph{drawdown periods} (market crashes), does the RL
          agent lose less than the benchmarks?
\end{itemize}

\begin{keyresult}
The RL agent may not always have the highest terminal wealth. The
goal is competitive returns with visibly better drawdown control.
Look for the RL curve to ``flatten'' during crashes rather than
plunging with the market.
\end{keyresult}

\subsection{Allocation Evolution Over Time}

\textbf{What it shows.} Stacked area chart of portfolio weights
$\bm{w}_t$ over time.

\textbf{What to look for.}
\begin{itemize}[leftmargin=2em]
    \item \textbf{Regime adaptation}: does the agent shift toward
          defensive assets (bonds, low-volatility stocks) before or during
          market stress?
    \item \textbf{Smoothness}: are the allocation changes gradual
          (indicating stable learned features) or erratic (indicating
          overfitting to noise)?
    \item \textbf{Comparison with benchmarks}: equal weight is flat by
          construction, risk parity is nearly flat. The RL agent should
          show meaningful variation correlated with market conditions.
\end{itemize}

\subsection{Risk--Return Scatter}

\textbf{What it shows.} Each strategy plotted with annualised return on
the $y$-axis and a risk metric (annualised volatility or CVaR) on the
$x$-axis.

\textbf{What to look for.}
\begin{itemize}[leftmargin=2em]
    \item A strategy in the \emph{upper-left} quadrant (high return, low
          risk) dominates.
    \item The RL agent should sit \emph{above} the line connecting
          equal weight and mean-variance, indicating a better
          risk--return trade-off.
    \item If using CVaR as the $x$-axis, the RL point should be near
          or below the threshold $\bar{c}$, confirming constraint
          satisfaction.
\end{itemize}

\subsection{Drawdown Comparison}

\textbf{What it shows.} Time series of drawdown
$\text{DD}_t = (V_{\max,t} - V_t) / V_{\max,t}$, where $V_t$ is
portfolio value and $V_{\max,t}$ is the running peak.

\textbf{What to look for.}
\begin{itemize}[leftmargin=2em]
    \item \textbf{Maximum drawdown}: the deepest trough for each
          strategy. The RL agent should have a shallower maximum drawdown.
    \item \textbf{Recovery time}: how long does it take to return to the
          previous peak? The RL agent should recover faster.
    \item \textbf{Frequency}: does the RL agent experience fewer
          significant drawdown events overall?
\end{itemize}

\begin{refbox}
For a comprehensive treatment of drawdown-based risk measures, see
Chekhlov, Uryasev \& Zabarankin~(2005), ``Drawdown Measure in Portfolio
Optimization.'' The maximum drawdown statistic is closely related to CVaR
--- both focus on tail events, but drawdown captures the temporal structure
of losses.
\end{refbox}


% ============================================================
\section{References}
\label{sec:references}
% ============================================================

\begin{enumerate}[leftmargin=2em, label={[\arabic*]}]

    \item \textbf{Markowitz, H.}~(1952).
          ``Portfolio Selection.''
          \emph{The Journal of Finance}, 7(1), 77--91.
          The foundational paper on mean-variance optimisation, establishing
          that investors should select portfolios based on both expected
          return and variance.

    \item \textbf{Merton, R.~C.}~(1969).
          ``Lifetime Portfolio Selection under Uncertainty: The
          Continuous-Time Case.''
          \emph{The Review of Economics and Statistics}, 51(3), 247--257.
          Extended Markowitz to a continuous-time dynamic setting, showing
          that optimal allocation depends on the state of the economy.

    \item \textbf{DeMiguel, V., Garlappi, L., \& Uppal, R.}~(2009).
          ``Optimal Versus Naive Diversification: How Inefficient Is the
          1/N Portfolio Strategy?''
          \emph{The Review of Financial Studies}, 22(5), 1915--1953.
          Demonstrated that equal-weight portfolios outperform many
          ``optimal'' strategies out-of-sample due to estimation error.

    \item \textbf{Tamar, A., Glassner, Y., \& Mannor, S.}~(2015).
          ``Optimizing the CVaR via Sampling.''
          \emph{Proceedings of the AAAI Conference on Artificial Intelligence},
          29(1).
          Developed algorithms for optimising CVaR-constrained MDPs using
          sampled trajectories, providing the theoretical basis for
          CVaR-constrained policy optimisation.

    \item \textbf{Salimans, T., Ho, J., Chen, X., Sidor, S., \&
          Sutskever, I.}~(2017).
          ``Evolution Strategies as a Scalable Alternative to Reinforcement
          Learning.''
          \emph{arXiv preprint arXiv:1703.03864}.
          Showed that NES can match or exceed the performance of
          gradient-based RL methods while being simpler to implement and
          easier to parallelise.

    \item \textbf{Wang, H., et al.}~(2025).
          ``ICVaR-DRL: Deep Reinforcement Learning with Iterated CVaR
          Constraints for Dynamic Portfolio Optimisation.''
          Extended CVaR-constrained RL with an iterated risk assessment
          framework, allowing the risk constraint to adapt over time.

\end{enumerate}

\end{document}
